\documentclass[
    10pt,
    a5paper,
    twoside,
    openright
]{memoir}

% ============================================================
% LIVRO DONA MARIA — PREÂMBULO
% ============================================================

\usepackage[T1]{fontenc}
\usepackage[utf8]{inputenc}
\usepackage[brazil]{babel}

\usepackage{libertinus}
\usepackage[protrusion=true,expansion=true]{microtype}

% ============================================================
% IDENTIDADE VISUAL
% ============================================================

\usepackage{xcolor}
\usepackage{tikz}

% Mundo — azul petróleo / azul acinzentado
\definecolor{CorMundo}{RGB}{67,91,101}

% Brasil — verde
\definecolor{CorBrasil}{RGB}{72,112,91}

% Dona Maria — terracota
\definecolor{CorMaria}{RGB}{154,96,78}

% Neutros
\definecolor{CorTexto}{RGB}{55,55,55}
\definecolor{CorCinza}{RGB}{120,120,120}

% Tema atual
\newcommand{\corTema}{CorMaria}

\newcommand{\temaMundo}{%
    \renewcommand{\corTema}{CorMundo}%
}

\newcommand{\temaBrasil}{%
    \renewcommand{\corTema}{CorBrasil}%
}

\newcommand{\temaMaria}{%
    \renewcommand{\corTema}{CorMaria}%
}

% Moldura temática dos capítulos
\newcommand{\molduraTematica}{%
    \begingroup
        \edef\x{\endgroup
            \noexpand\AddToHookNext{shipout/background}{%
                \noexpand\begin{tikzpicture}[remember picture,overlay]
                    \noexpand\draw[
                        color=\corTema,
                        line width=0.45pt
                    ]
                    ([xshift=11mm,yshift=-11mm]current page.north west)
                    rectangle
                    ([xshift=-11mm,yshift=11mm]current page.south east);
                \noexpand\end{tikzpicture}%
            }%
        }%
    \x
}

% Moldura neutra das partes
\newcommand{\molduraParte}{%
    \begingroup
        \edef\x{\endgroup
            \noexpand\AddToHookNext{shipout/background}{%
                \noexpand\begin{tikzpicture}[remember picture,overlay]
                    \noexpand\draw[
                        color=CorCinza,
                        line width=0.35pt
                    ]
                    ([xshift=11mm,yshift=-11mm]current page.north west)
                    rectangle
                    ([xshift=-11mm,yshift=11mm]current page.south east);
                \noexpand\end{tikzpicture}%
            }%
        }%
    \x
}

% ============================================================
% GEOMETRIA, TEXTO E PAGINAÇÃO
% ============================================================

\setstocksize{210mm}{148mm}
\settrimmedsize{210mm}{148mm}{*}

\setlrmarginsandblock{24mm}{19mm}{*}
\setulmarginsandblock{22mm}{22mm}{*}

\setheadfoot{12mm}{10mm}
\setheaderspaces{*}{8mm}{*}

\checkandfixthelayout

\setlength{\parindent}{1.2em}
\setlength{\parskip}{0pt}
\setlength{\emergencystretch}{2em}

\linespread{1.08}

\makepagestyle{livro}

\makeevenhead{livro}
    {\small\thepage}
    {}
    {\small\itshape\leftmark}

\makeoddhead{livro}
    {\small\itshape\rightmark}
    {}
    {\small\thepage}

\makeevenfoot{livro}{}{}{}
\makeoddfoot{livro}{}{}{}

\pagestyle{livro}

\aliaspagestyle{chapter}{empty}
\aliaspagestyle{part}{empty}

% ============================================================
% ESTILO DE CAPÍTULOS, PARTES E SUBTÍTULOS
% ============================================================

\makechapterstyle{dona-maria}{%
    \renewcommand{\printchaptername}{}
    \renewcommand{\printchapternum}{}
    \renewcommand{\afterchapternum}{}

    \renewcommand{\chaptitlefont}{%
        \normalfont
        \fontsize{22}{26}\selectfont
        \bfseries
        \color{\corTema}
    }

    \renewcommand{\afterchaptertitle}{%
        \molduraTematica
        \par
        \vskip 8pt
        \noindent
        {\color{\corTema}\rule{\textwidth}{0.8pt}}
        \vskip 22pt
    }
}
\chapterstyle{dona-maria}

\renewcommand{\partnamefont}{%
    \normalfont
    \small
    \scshape
    \color{CorCinza}
}

\renewcommand{\partnumfont}{%
    \normalfont
    \fontsize{16}{20}\selectfont
    \color{CorTexto}
}

\renewcommand{\parttitlefont}{%
    \molduraParte
    \normalfont
    \fontsize{30}{34}\selectfont
    \bfseries
    \color{CorTexto}
}

\setsecnumdepth{subsection}

\setsecheadstyle{%
    \normalfont
    \large
    \bfseries
    \color{\corTema}
}

\setsubsecheadstyle{%
    \normalfont
    \normalsize
    \bfseries
    \color{\corTema}
}

% ============================================================
% IMAGENS E LEGENDAS
% ============================================================

\usepackage{graphicx}
\graphicspath{{imagens/}}

\usepackage[
    font=small,
    labelfont=bf,
    justification=centering
]{caption}

\captionsetup{skip=8pt}

% Fotografia grande
\newcommand{\fotoGrande}[3]{%
    \begin{figure}[htbp]
        \centering
        \includegraphics[width=0.88\textwidth]{#1}
        \caption{#2}
        \label{#3}
    \end{figure}
}

% Fotografia média
\newcommand{\fotoMedia}[3]{%
    \begin{figure}[htbp]
        \centering
        \includegraphics[width=0.68\textwidth]{#1}
        \caption{#2}
        \label{#3}
    \end{figure}
}

% Fotografia pequena
\newcommand{\fotoPequena}[3]{%
    \begin{figure}[htbp]
        \centering
        \includegraphics[width=0.48\textwidth]{#1}
        \caption{#2}
        \label{#3}
    \end{figure}
}

% Fonte de imagem
\newcommand{\fonteImagem}[1]{%
    \par\smallskip
    {\footnotesize\itshape Fonte: #1.}%
}

% ============================================================
% REFERÊNCIAS E NOTAS
% ============================================================

\usepackage{csquotes}

% Bibliografia centralizada em .bib.
% O estilo authoryear-icomp mantém as referências discretas.
\usepackage[
    backend=biber,
    style=authoryear-icomp,
    autocite=footnote,
    giveninits=true,
    maxbibnames=50,
    maxcitenames=2,
    language=portuguese
]{biblatex}

\addbibresource{bibliografia/fontes.bib}

\SetBlockEnvironment{quotation}

\renewenvironment{quotation}
    {\list{}{
        \setlength{\leftmargin}{1.5em}
        \setlength{\rightmargin}{1.5em}
    }
    \item\relax
    \small
    \itshape}
    {\endlist}

% Notas de rodapé discretas
\setlength{\footmarkwidth}{1.5em}
\setlength{\footmarksep}{0.3em}
\renewcommand{\foottextfont}{\footnotesize}

% ============================================================
% TABELAS, LINKS, PDF E COMANDOS EDITORIAIS
% ============================================================

\usepackage{booktabs}
\usepackage{array}

\usepackage[
    hidelinks,
    pdfusetitle
]{hyperref}

\usepackage{bookmark}

\hypersetup{
    pdftitle={Dona Maria},
    pdfauthor={Autor},
    pdfsubject={A história de Dona Maria, 1950--2025},
    pdfkeywords={Dona Maria, memória, família, Brasil, história}
}

\newcommand{\ano}[1]{\textsc{#1}}

\newcommand{\linhaeditorial}{%
    \par\vspace{1em}\noindent
    {\color{\corTema}\rule{\textwidth}{0.5pt}}
    \par\vspace{1em}
}

\newcommand{\memoria}{%
    \begin{center}
        {\color{\corTema}\small$\ast\quad\ast\quad\ast$}
    \end{center}
}


% ============================================================
% FIM DO PREÂMBULO
% ============================================================


\begin{document}

\frontmatter

\chapter*{Dona Maria}
\addcontentsline{toc}{chapter}{Dona Maria}

Uma vida entre 1950 e 2025.

\tableofcontents*

\mainmatter

\part{1950--1969}

\temaMundo
\chapter*{Do pós-guerra à lua}
\addcontentsline{toc}{chapter}{Do pós-guerra à lua}

Na aurora dos anos 1950, o mundo ainda se recuperava das cicatrizes deixadas pela Segunda Guerra Mundial. O conflito havia redesenhado fronteiras e abalado nações, deixando um rastro de destruição que ecoava no cotidiano de milhões de pessoas. As cidades europeias eram um espelho partido, onde a devastação parecia impregnar cada ruína, cada parede quebrada. Mas o espírito humano, resiliente como nunca, encontrava forças para sonhar com a reconstrução e com um futuro diferente daquele que o conflito mundial havia pintado.

Logo após o fim da guerra, em 1945, nasceu a Organização das Nações Unidas (ONU), carregando o pesado fardo de garantir a paz global. Inspirada nas falhas da Liga das Nações, a ONU foi concebida como uma arena onde o diálogo substituiria o confronto, e as diplomacias emergiriam vitoriosas sobre as armas. Reuniões entre líderes mundiais buscavam consolidar o equilíbrio geopolítico, enquanto a Cortina de Ferro começava a se abaixar entre o Oriente e o Ocidente, dando início à Guerra Fria.

Na Palestina, outro conflito surgia. Em 1948, a fundação do Estado de Israel, respaldada por uma votação da ONU, transformou a região em um campo de tensões entre árabes e judeus. Os confrontos que se seguiram marcaram o início de uma longa e sangrenta disputa territorial, que ainda reverberaria por décadas. O novo mapa político global estava repleto de tensões mal resolvidas, onde cada decisão parecia ecoar novos conflitos.

Enquanto isso, as nações europeias, com o apoio do Plano Marshall, se reerguiam das cinzas. A Alemanha, outrora dividida, começava sua lenta reconstrução, e cidades como Londres e Paris ressurgiam como faróis culturais e econômicos. A Europa, embora dilacerada pelo conflito, reencontrava sua identidade e reconstruía suas economias, agora sob a sombra protetora do capitalismo ocidental e o medo crescente da expansão comunista.

Simultaneamente, no hemisfério Sul, movimentos de independência varriam o continente africano e asiático. Colônias, antes presas às mãos dos impérios europeus, começavam a lutar por sua autonomia. Na Índia, em 1947, a liderança de Mahatma Gandhi conquistou a liberdade, inspirando outras nações a seguir o mesmo caminho. Durante os anos 1950 e 1960, países como o Gana, Congo, Argélia e Indonésia se libertaram de seus colonizadores, embora a transição para a independência fosse muitas vezes marcada por conflitos internos e guerras civis, como foi o caso da Guerra do Vietnã, que se estenderia até os anos 1970.

A corrida espacial surgiu como a disputa simbólica mais brilhante da Guerra Fria. Estados Unidos e União Soviética competiam não apenas pela supremacia militar e ideológica, mas também pelo domínio dos céus e além. Em 1957, o satélite soviético Sputnik cruzou o céu, assombrando o Ocidente e acendendo o fogo da exploração espacial. No início dos anos 1960, o cosmonauta Yuri Gagarin tornou-se o primeiro humano a orbitar a Terra, um triunfo soviético que fez o mundo prender a respiração. Mas os Estados Unidos responderam à altura. Sob a presidência de John F. Kennedy, o projeto Apollo foi lançado, e em 1969, a humanidade deu seu passo mais ousado: Neil Armstrong pisou na Lua, marcando o ápice da corrida espacial.

Entretanto, enquanto o olhar do mundo se voltava para as estrelas, conflitos terrenos se intensificavam. As guerras por procuração moldavam o tabuleiro global. Na Coreia, entre 1950 e 1953, as tropas americanas e soviéticas se enfrentaram indiretamente, consolidando a divisão da península em duas nações antagônicas. No Vietnã, o conflito entre o Norte comunista e o Sul capitalista se intensificou, arrastando os Estados Unidos para uma guerra impopular e devastadora.

O Oriente Médio e a América Latina também tornaram-se palcos de conflitos indiretos entre as superpotências. Revoluções e golpes de estado, muitas vezes patrocinados pela CIA ou pela KGB, transformaram o Chile, Cuba, Egito e outros países em verdadeiros campos de batalha ideológicos.

Os anos 1950 e 1960 foram marcados por uma dualidade paradoxal. Enquanto a ciência e a tecnologia empurravam os limites do que a humanidade poderia alcançar, o planeta se afundava em novas formas de conflito e opressão. A Guerra Fria, com suas promessas de destruição mútua, mantinha o mundo em suspenso. Mas, ao mesmo tempo, a ideia de que o futuro poderia ser algo maior e mais grandioso começou a tomar forma, impulsionada por avanços científicos e pelo sonho de explorar novos mundos.

Esse período foi uma dança constante entre esperança e medo, progresso e destruição, onde o homem parecia capaz de moldar o destino da Terra ou destruí-la completamente. O fim dos anos 1960 se aproximava com uma humanidade mais ousada e, ao mesmo tempo, mais frágil, prestes a redefinir o curso de sua própria história.
\temaBrasil
\chapter*{Do petróleo à ditadura}
\addcontentsline{toc}{chapter}{Do petróleo à ditadura}

Este é um segundo texto de teste.

\temaMaria
\chapter*{Do nascimento à adolescência}
\addcontentsline{toc}{chapter}{Do nascimento à adolescência}

Aqui começará a história de Dona Maria.

\part{1970--1989}

\temaMundo
\chapter*{Conflitos e quedas}
\addcontentsline{toc}{chapter}{Conflitos e quedas}

Texto de teste.

\temaBrasil
\chapter*{Do tri à eleição direta}
\addcontentsline{toc}{chapter}{Do tri à eleição direta}

Texto de teste.

\temaMaria
\chapter*{Casamento, filhos e separação}
\addcontentsline{toc}{chapter}{Casamento, filhos e separação}

Texto de teste.

\part{1990--2009}

\temaMundo
\chapter*{Da era da informação às crises econômicas}
\addcontentsline{toc}{chapter}{Da era da informação às crises econômicas}

Texto de teste.

\temaBrasil
\chapter*{Do impeachment à marolinha}
\addcontentsline{toc}{chapter}{Do impeachment à marolinha}  

Texto de teste.

\temaMaria
\chapter*{Cirurgias, derrotas e conquistas}
\addcontentsline{toc}{chapter}{Cirurgias, derrotas e conquistas}

Texto de teste.

\part{2010--2025}

\temaMundo
\chapter*{Guerras e mudanças climáticas}
\addcontentsline{toc}{chapter}{Guerras e mudanças climáticas}

Texto de teste.

\temaBrasil
\chapter*{Manifestações, impeachment e extrema direita}
\addcontentsline{toc}{chapter}{Manifestações, impeachment e extrema direita}

Texto de teste.

\temaMaria
\chapter*{Angústias, terapia e recomeço}
\addcontentsline{toc}{chapter}{Angústias, terapia e recomeço}

Texto de teste.

\end{document}